\section*{Hoofdstuk \ref{ch:b1:replication-mitosis} --- DNA-replicatie en mitose}

\begin{solution}{exo:b1:replication-mitosis:1}
Polymerasen voegen \omterm{def:b1:nucleic-acids:nucleotide}{nucleotiden} alleen aan een $3'$-uiteinde toe, zodat nieuwe
strengen $5' \to 3'$ groeien; de twee matrijzen zijn antiparallel, zodat bij een
vork de ene nieuwe streng ononderbroken naar de vork toe kan groeien terwijl de
andere er juist van weg moet groeien en telkens opnieuw in fragmenten begint
naarmate er meer matrijs vrijkomt.
\end{solution}

\begin{solution}{exo:b1:replication-mitosis:2}
Helicase (ontwindt), enkelstrengbindende \omterm{def:b1:proteins:peptide}{eiwitten} (houden de strengen uit elkaar),
topo-isomerase (heft de torsie vooruit op), primase (RNA-primers), \omterm{def:b1:replication-mitosis:fork}{DNA-polymerase}
(verlengt primers, leest proef), een tweede polymerase (vervangt de primers),
ligase (verbindt de fragmenten).
\end{solution}

\begin{solution}{exo:b1:replication-mitosis:3}
G$_1$: $q$, 2n \omterm{def:b1:genomes:genome}{chromosomen} van één \omterm{def:b1:replication-mitosis:mitosis}{chromatide}. G$_2$: $2q$, 2n \omterm{def:b1:genomes:genome}{chromosomen} van
twee \omterm{def:b1:replication-mitosis:mitosis}{chromatiden}. Na de \omterm{def:b1:replication-mitosis:mitosis}{mitose}: $q$ per dochtercel, elk 2n \omterm{def:b1:genomes:genome}{chromosomen}.
\end{solution}

\begin{solution}{exo:b1:replication-mitosis:4}
Profase (condensatie, de \omterm{def:b1:replication-mitosis:mitosis}{spoelfiguur} vormt zich), prometafase (de envelop valt
uiteen, de kinetochoren hechten), metafase (uitlijning op de evenaar), anafase (de
\omterm{def:b1:replication-mitosis:mitosis}{chromatiden} gaan uiteen), telofase (de enveloppen vormen zich opnieuw), daarna de
cytokinese.
\end{solution}

\begin{solution}{exo:b1:replication-mitosis:5}
De helft van het \omterm{def:b1:genomes:genome}{genoom} is synthese van de volgende streng: $\num{6.4e9}/150 =
\num{4.3e7}$ fragmenten; evenveel primers, verwijderingen en ligaties.
\end{solution}

\begin{solution}{exo:b1:replication-mitosis:6}
$\num{6.4e9}\times 10^{-5} = \num{64000}$; $\times 10^{-7}$: 640;
$\times 10^{-9}$: ongeveer 6 mutaties per deling.
\end{solution}

\begin{solution}{exo:b1:replication-mitosis:7}
$5000/80 = 62$ delingen. Een tumor moet onbeperkt delen; zonder telomerase zouden
zijn telomeren opraken en zouden zijn \omterm{def:b1:cell-unit-of-life:cell}{cellen} stilvallen of sterven, en daarom
activeren vrijwel alle kankers het \omterm{def:b1:enzymes:enzyme}{enzym} opnieuw.
\end{solution}

\begin{solution}{exo:b1:replication-mitosis:8}
Cyclus $= 1\,\mathrm{h}/0.05 = \qty{20}{h}$; S $= 0.30\times 20 =
\qty{6}{h}$.
\end{solution}

\begin{solution}{exo:b1:replication-mitosis:9}
G$_1$: 8 \omterm{def:b1:genomes:genome}{chromosomen}, 8 \omterm{def:b1:replication-mitosis:mitosis}{chromatiden}, $q$. G$_2$: 8, 16, $2q$. Metafase: 8, 16,
$2q$. Dochtercel: 8, 8, $q$.
\end{solution}

\begin{solution}{exo:b1:replication-mitosis:10}
Er beginnen elke \qty{20}{min} nieuwe rondes bij de oorsprong voordat de vorige
ronde is afgelopen: de replicatie overlapt, en een pasgeboren \omterm{def:b1:cell-unit-of-life:cell}{cel} heeft al
gedeeltelijk gerepliceerd \omterm{def:b1:nucleic-acids:chain}{DNA}. Bij de deling is het \omterm{def:b1:genomes:genome}{chromosoom} dat zijn ronde
afmaakt \qty{40}{min} eerder begonnen, heeft een tweede ronde die \qty{20}{min}
eerder begon haar vorken halverwege, en begint een derde net: 4 oorsprongen en 6
vorken in de delende \omterm{def:b1:cell-unit-of-life:cell}{cel}.
\end{solution}

\begin{solution}{exo:b1:replication-mitosis:11}
Er vormt zich geen \omterm{def:b1:replication-mitosis:mitosis}{spoelfiguur}: de \omterm{def:b1:genomes:genome}{chromosomen} condenseren wel maar kunnen zich
niet uitlijnen of scheiden, en de \omterm{def:b1:cell-unit-of-life:cell}{cel} blijft bij het spoelfiguurcontrolepunt steken
in een metafase-achtige toestand. Glijdt zij zonder te delen terug in de \omterm{def:b1:replication-mitosis:cycle}{interfase},
dan heeft zij 4n \omterm{def:b1:genomes:genome}{chromosomen} (tetraploïd). Stilgelegde \omterm{def:b1:cell-unit-of-life:cell}{cellen} met gecondenseerde,
afzonderlijke \omterm{def:b1:genomes:genome}{chromosomen} zijn precies wat een \omterm{def:b1:genomes:karyotype}{karyotype} nodig heeft: het middel
wordt gebruikt om ze op te hopen.
\end{solution}

\begin{solution}{exo:b1:replication-mitosis:12}
De replicatie waarborgt dat elk \omterm{def:b1:genomes:genome}{chromosoom} twee identieke \omterm{def:b1:replication-mitosis:mitosis}{chromatiden} wordt, tot op
één fout per miljard; de \omterm{def:b1:replication-mitosis:mitosis}{spoelfiguur} waarborgt dat elke dochtercel er van elk paar
één ontvangt, door elk \omterm{def:b1:genomes:genome}{chromosoom} vanuit beide polen vast te hechten en de anafase
op te houden tot elk \omterm{def:b1:genomes:genome}{chromosoom} onder spanning staat. Het \omterm{prop:b1:replication-mitosis:checkpoints}{controlepunt} aan het eind
van G$_2$ zorgt dat het kopiëren klaar is voordat het tellen begint; het
spoelfiguurcontrolepunt zorgt dat het tellen is opgezet voordat de scheiding komt.
Een \omterm{def:b1:cell-unit-of-life:cell}{cel} die verkeerd verdeelt eindigt met één \omterm{def:b1:genomes:genome}{chromosoom} te veel of te weinig
(aneuploïdie): een tot op de letter juiste kopie, op het verkeerde adres bezorgd
--- wat de meeste van zulke \omterm{def:b1:cell-unit-of-life:cell}{cellen} doodt en de meeste kankers kenmerkt.
\end{solution}

\begin{solution}{pb:b1:replication-mitosis:1}
\textbf{1.} $50\times 28\,800 = \qty{1.44e6}{bp}$.
\textbf{2.} \qty{2.88e6}{bp}.
\textbf{3.} $\num{6.4e9}/\num{2.88e6} = 2220$ oorsprongen.
\textbf{4.} \qty{2.9}{Mb}, oftewel $2.9\times 10^6\times 0.34\,\mathrm{nm} =
\qty{0.98}{mm}$ \omterm{def:b1:nucleic-acids:chain}{DNA} tussen de oorsprongen.
\textbf{5.} $\num{6.4e9}/\num{40000} = \qty{160}{kb}$ afstand; elke vork legt
ongeveer \qty{80}{kb} af, een half uur werk.
\textbf{6.} \qty{125}{Mb} per vork bij \qty{50}{bp/s}: $\num{2.5e6}$ s, 29 dagen.
\textbf{7.} De vorksnelheid wordt begrensd door de chemie en door het chromatine
(twintigvoudig trager dan bij bacteriën); oorsprongen kunnen onbeperkt worden
vermenigvuldigd, zodat de S-fase wordt verkort door startpunten toe te voegen en
niet door de vorken op te jagen.
\textbf{8.} De synthese van de volgende streng beslaat de helft van het \omterm{def:b1:genomes:genome}{genoom},
$\num{3.2e9}$ \omterm{def:b1:nucleic-acids:nucleotide}{nucleotiden}: $\num{2.1e7}$ fragmenten (de twee vorken van elke
oorsprong hebben elk een volgende streng, en de helften tellen op tot één \omterm{def:b1:genomes:genome}{genoom}).
\textbf{9.} $\num{2.1e7} + 2\times\num{40000} = \num{2.1e7}$ primers.
\textbf{10.} Ongeveer $\num{2.1e7}$ ligaties (één per fragment) plus één per
ontmoeting van vorken.
\textbf{11.} $2\times\num{6.4e9} = \num{1.28e10}$ \omterm{def:b1:nucleic-acids:nucleotide}{nucleotiden}.
\textbf{12.} $\num{2.56e10}$ \omterm{def:b1:water-small-molecules:atp}{ATP}; over \qty{28800}{s} is dat $\num{8.9e5}$ per
seconde.
\textbf{13.} Ongeveer \qty{0.1}{\%} van de ATP-omzet van de \omterm{def:b1:cell-unit-of-life:cell}{cel}: het polymeriseren
is goedkoop; duur zijn het maken van de \omterm{def:b1:nucleic-acids:nucleotide}{nucleotiden} en van de \omterm{def:b1:genomes:nucleosome}{histonen}.
\textbf{14.} \num{64000}, 640 en 6,4 fouten.
\textbf{15.} $1500\times 10^{-9} = \num{1.5e-6}$: één \omterm{def:b1:genomes:gene}{gen} op zevenhonderdduizend
per deling.
\textbf{16.} $10^{16}\times 6.4 = \num{6.4e16}$ mutaties --- elke mogelijke
enkelvoudige verandering van het \omterm{def:b1:genomes:genome}{genoom} vele malen, verspreid over $10^{13}$
\omterm{def:b1:cell-unit-of-life:cell}{cellen}: vrijwel alle vallen in niet-coderend \omterm{def:b1:nucleic-acids:chain}{DNA} of in \omterm{def:b1:cell-unit-of-life:cell}{cellen} die worden
afgestoten, en een gemuteerde \omterm{def:b1:cell-unit-of-life:cell}{cel} is er één op miljarden; alleen enkele combinaties
in één \omterm{def:b1:cell-unit-of-life:cell}{cel} doen ertoe.
\textbf{17.} Honderdvoudig ($10^{-7}$): de reeks mutaties die een kanker maakt, en
die bij $10^{-9}$ decennia kost om zich op te hopen, hoopt zich veel eerder op ---
erfelijke gebreken in het mismatchherstel veroorzaken darmkanker op jonge
volwassen leeftijd.
\textbf{18.} $\num{2.3e6}$ paren per vork bij \qty{1000}{bp/s}: \qty{2300}{s},
\qty{38}{min}.
\textbf{19.} Ongeveer twee overlappende rondes; een pasgeboren \omterm{def:b1:cell-unit-of-life:cell}{cel} draagt twee
oorsprongen (die elk al eenmaal zijn gerepliceerd) en heeft haar vorken halverwege:
haar \omterm{def:b1:genomes:genome}{chromosoom} is bij de geboorte al gedeeltelijk gerepliceerd.
\textbf{20.} $\num{4.6e6}/1500 = 3070$ fragmenten per ronde, ongeveer 1,3 per
seconde.
\textbf{21.} $\num{4.6e6}\times 10^{-9} = \num{4.6e-3}$ fouten per ronde: ongeveer
één dochtercel op 200 draagt een nieuwe mutatie.
\textbf{22.} $10^9\times\num{4.6e-3} = \num{4.6e6}$ nieuwe mutaties per
milliliter; een bepaald \omterm{def:b1:genomes:gene}{gen} van \qty{1}{kb} wordt $\num{4.6e6}\times
1000/\num{4.6e6} = 1000$ keer getroffen.
\textbf{23.} Met duizend onafhankelijke mutaties in elk willekeurig \omterm{def:b1:genomes:gene}{gen} per
milliliter per deling is elke mogelijke enkelvoudige verandering --- ook die welke
resistentie geven --- al aanwezig in een grote kweek voordat het antibioticum wordt
toegediend; het middel selecteert, het schept niet.
\textbf{24.} Mens: $\num{6.4e9}$ paren, \num{40000} oorsprongen, \qty{50}{bp/s},
\qty{8}{h}. Bacterie: $\num{4.6e6}$ paren, één oorsprong, \qty{1000}{bp/s},
\qty{40}{min}. Een \omterm{def:b1:genomes:genome}{genoom} dat duizend keer groter is, gekopieerd door vorken die
twintig keer trager zijn, in slechts twaalf keer de tijd: het verschil zijn de
tienduizenden oorsprongen die parallel werken.
\textbf{25.} Minimaal ongeveer \num{2200} oorsprongen, als zij allemaal tegelijk
zouden starten; in werkelijkheid zo'n \num{40000}, één per \qty{160}{kb}, die na
elkaar starten zodat elke vork ongeveer \qty{80}{kb} aflegt.
\end{solution}
